% !TeX root = ../main.tex
\begin{englishabstract}
% Uncomment if funding information is available.
% \footnotetext{*The work was supported by the Foundation (foundation ID).}
% \astfootnote{The work was supported by the Foundation (foundation ID).}

With the widespread deployment of deep neural networks and large language models in the cloud and at the edge, the demand of artificial intelligence for hardware compute is growing rapidly. Neural network chips have been proposed as domain-specific architectures and are evolving from single-core to multi-core designs. However, evaluating end-to-end performance of a multi-core chip before tape-out faces a dilemma: hardware emulation platforms offer high accuracy but are costly, slow to iterate, and difficult to reconfigure flexibly; analytical models offer low iteration cost yet fail to capture system-level dynamic behavior such as on-chip network congestion, off-chip memory contention, and multi-core synchronization. Cycle-level simulators sit between these two extremes and serve as critical infrastructure for performance evaluation prior to tape-out. Yet, open-source simulators that target multi-core neural network chips and directly consume scheduler intermediate representations remain absent.\newline

To fill this gap, this thesis designs and implements a cycle-level simulator for multi-core neural network chips driven by a scheduler intermediate representation. Through a unified task abstraction, the simulator casts multi-core execution as a model of core-level workload, explicit dependency, and data flow, and advances compute, interconnect, and memory in a coordinated manner under a unified timeline with multi-clock synchronization. The simulator records per-core cycles in loading, computing, writeback, and network backpressure, outputs chip-state statistics, and visualizes simulation results as heatmaps and Gantt charts, forming a closed loop from parsing to execution to statistics, so that bottleneck attribution no longer relies on human experience.\newline

Experiments thoroughly validate the accuracy of the simulator and its value for design-space exploration. For accuracy validation, execution results from the commercial Synopsys ZeBu hardware emulation platform are taken as the reference data. Across different workloads and different batch-size configurations, the total-cycle deviation between the simulator output and the reference data stays within $\pm$5.2\%, reaching accuracy close to that of hardware emulation. Building on this, this thesis conducts systematic bottleneck analysis and design-space exploration for multi-core neural network chips. On a four-core test chip running the ResNet50 model, a progressive optimization path that first upgrades off-chip memory bandwidth and then widens the on-chip memory port is simulated, reducing end-to-end total cycles by 73.7\% (a 3.80-fold speedup) and revealing a stepwise bottleneck migration from off-chip bandwidth to on-chip port bandwidth. A further design-space exploration is carried out along three dimensions---single-core compute capability, core-count scaling, and batch size---with 12 configurations, quantitatively yielding three findings: (1) when memory bandwidth falls below what compute intensity demands, increasing compute capability contributes nothing to throughput; (2) once chip core count exceeds the workload's parallelism limit, performance degrades severely and may even deadlock; (3) limited by memory-bandwidth contention, the batch-scaling efficiency of the workload is only about 24\%. In summary, this thesis provides a reproducible cycle-level simulation foundation and practical evaluation capability for system-level bottleneck localization and software--hardware co-optimization of multi-core neural network chips prior to tape-out.

\englishKeywords{Neural network chip; Multi-core architecture; Cycle-level simulator; Design space exploration}

% 每个关键词组的第一个字母大写，其余为小写，每一关键词之间用分号分开，最后一个关键词后不打标点符号。例如：Drip irrigation emitter; RP\&M; Hydraulics; Labyrinth flow channel\newline

\englishType{Application Research}

% 须与中文摘要中的论文类型一致；每个单词第一个字母大写，其余为小写。例如：Applied Research

% 论文类型包括：a.理论研究(Theoretical Research)；b.应用基础(Application Fundamentals)；c.应用研究(Application Research)；d.研究报告(Research Report)；e.设计报告(Design Report)；f.案例分析(Case Study)；g.调研报告(Investigation Report)；h.产品研发(Product Development)；i.工程设计(Engineering Design)；j.工程/项目管理(Engineering/Project Management)；k.其它（Others）。

\end{englishabstract}
